\documentclass[12pt,a4paper]{article}
\usepackage[utf8]{inputenc}
\usepackage{geometry}
\usepackage{graphicx}
\usepackage{listings}
\usepackage{xcolor}
\usepackage{float}
\usepackage{hyperref}
\hypersetup{
    colorlinks=true,
    linkcolor=black,
    urlcolor=blue,
    citecolor=blue
}
\geometry{margin=1in}

\setcounter{secnumdepth}{2}
\setcounter{tocdepth}{2}

\renewcommand{\thesubsection}{\arabic{section}.\arabic{subsection}}

\usepackage{etoolbox}
\pretocmd{\appendix}{%
  \renewcommand{\thesection}{\arabic{section}}%
}{}{}

\definecolor{codegreen}{rgb}{0,0.6,0}
\definecolor{codegray}{rgb}{0.5,0.5,0.5}
\definecolor{codepurple}{rgb}{0.58,0,0.82}
\definecolor{backcolour}{rgb}{0.95,0.95,0.92}

\lstdefinestyle{mystyle}{
    backgroundcolor=\color{backcolour},   
    commentstyle=\color{codegreen},
    keywordstyle=\color{magenta},
    numberstyle=\tiny\color{codegray},
    stringstyle=\color{codepurple},
    basicstyle=\ttfamily\footnotesize,
    breaklines=true,
    numbers=left,
    numbersep=5pt,
    showstringspaces=false
}

\lstset{style=mystyle}

\begin{document}

\begin{titlepage}
\centering

\includegraphics[width=0.25\textwidth]{iut_logo.png}\\[1cm]

{\Large \textbf{Islamic University of Technology (IUT)}}\\[0.3cm]
{\large Department of Computer Science and Engineering}\\[1.2cm]

{\large \textbf{CSE 4802: Compiler Design Lab}}\\[0.5cm]

{\LARGE \textbf{LAB 4 REPORT}}\\[0.3cm]

{\Large YACC-Based Arithmetic Expression Syntax Analyzer with Symbol Table Management}\\[1.5cm]

\textbf{Submitted by:}\\[0.3cm]
MD. Nazmus Sadiq\\
Student ID: 210041139\\
Section: 1A\\[1.5cm]

BSc. in Computer Science and Engineering\\[1.5cm]

{\large June 19, 2026}

\vfill
\end{titlepage}

\tableofcontents
\newpage

% ================= INTRODUCTION =================

\section{Introduction}

This lab extends the basic YACC-based arithmetic expression syntax analyzer by introducing variable support and a dynamic symbol table. In addition to parsing and evaluating arithmetic expressions, the system now supports assignment statements, identifiers, and persistent storage of variables. The symbol table is implemented using a linked list, allowing dynamic insertion and retrieval of variables at runtime. Flex is used for lexical analysis, while YACC handles syntax parsing and semantic evaluation.

% ================= PROBLEM =================

\section{Problem Statement}

The objective of this lab is to extend a YACC-based arithmetic expression parser with the following features:
\begin{itemize}
    \item Identify INTEGER, DOUBLE, and IDENTIFIER tokens.
    \item Support assignment statements (e.g., marks1 = 20.5, total = marks1 + marks2).
    \item Implement dynamic symbol table management using a linked list.
    \item Save symbol table contents to a file (symbols.txt) after each update.
    \item Load symbol table contents when the parser starts.
\end{itemize}

% ================= BACKGROUND =================

\section{Background Theory}

A syntax analyzer (parser) verifies whether a sequence of tokens follows grammatical rules. In compiler design, Flex performs lexical analysis by converting input into tokens, and YACC constructs a parser based on grammar rules.

In this extended implementation, a symbol table is introduced to store identifiers and their values during execution. The symbol table is implemented using a linked list to allow dynamic memory allocation. Each variable is stored with its name, value, and type (INTEGER or DOUBLE). Additionally, file handling is used to persist the symbol table between executions using a text file (symbols.txt).

% ================= APPENDIX =================

\appendix

\section{Lex Source Code (lexer.l)}
\lstinputlisting[language=C++]{lexer.l}

\section{YACC Source Code (parser.y)}
\lstinputlisting[language=C++]{parser.y}

\section{Sample Input (input.txt)}
\lstinputlisting[language=C++]{input.txt}

\section{Output (output.txt)}
\lstinputlisting[language=C++]{output.txt}

\section{Symbol Table File (symbols.txt)}
\lstinputlisting[language=C++]{symbols.txt}

\end{document}